\RequirePackage[T1]{fontenc}
\documentclass[conference]{IEEEtran}
\usepackage{cite,amsmath,amssymb,graphicx,booktabs,url,balance}
\usepackage[hidelinks]{hyperref}
\hyphenation{classi-fi-ca-tion cali-bra-tion}
\graphicspath{{../docs/phase-2/evaluation/}{../docs/phase-2/ablations/}{../docs/phase-2/}}
\begin{document}
\title{Smart Waste Segregation Assistant: Phase~2\\Deep Transfer Learning, Hybrid Fusion, and Frozen Evaluation}
\author{\IEEEauthorblockN{Om Yadav, Aditya Bharadwaj, Shourya Bafna,\\Daksh Batra, Prajjwal Tripathi, Saksham Miglani}
\IEEEauthorblockA{CSAI Project 12 --- Phase 2: Deep Learning and Evaluation\\13 November 2026}}
\maketitle
\pagestyle{plain}
\thispagestyle{plain}
\begin{abstract}
This paper reports Phase~2 of a six-class waste material classifier on the TrashNet benchmark. Building on Phase~1's RBF SVM (test macro F1 0.8963), we evaluate deep transfer learning with MobileNetV2 and EfficientNet-B0, a systematic one-factor-at-a-time ablation study, and a late-fusion hybrid combining the fine-tuned CNN with the Phase~1 SVM. Temperature scaling is applied on training out-of-fold (OOF) predictions, and two abstention thresholds are pre-registered on the validation set before a single frozen test pass. The H3 late-fusion hybrid achieves macro F1 \textbf{0.9233} [0.889, 0.952] and accuracy \textbf{93.63\%} on 377 held-out test images, significantly outperforming the A2.1 CNN ($p=0.000019$, McNemar exact). The advantage over the Phase~1 SVM ($p=0.035$) is suggestive but does not survive Bonferroni correction at $\alpha_{Bonf}=0.0167$. An educational Gradio demo with per-class probabilities and an uncertainty banner uses frozen calibrated thresholds from \texttt{protocol.json}.
\end{abstract}
\begin{IEEEkeywords}
Waste classification, MobileNetV2, EfficientNet, transfer learning, late fusion, temperature scaling, selective classification, McNemar test.
\end{IEEEkeywords}

%% ======================================================
\section{Introduction}
%% ======================================================
\label{sec:intro}
Automated visual waste sorting can support municipal recycling by routing items to the correct stream. Phase~1 established a reproducible evaluation protocol on the TrashNet benchmark~\cite{yang2016trashnet}: a group-aware 70/15/15 split of 2,521 images, five frozen training folds, and an RBF SVM on 1,608 handcrafted color and texture features that reached test macro F1 0.8999. Phase~2 investigates whether deep transfer learning and hybrid fusion can improve on this baseline while respecting the same frozen test split and pre-registration discipline.

Three questions structure Phase~2: (1)~Can a pretrained CNN adapted for TrashNet beat handcrafted features? (2)~Does late probability fusion correct the residual errors of either modality? (3)~Does temperature calibration on training OOF predictions transfer reliably to test? We answer each via a pre-registered protocol~\cite{elyaniv2010selective} in which all configuration decisions are made before the test split is touched.

%% ======================================================
\section{Related Work}
%% ======================================================
\label{sec:related}
MobileNetV2 uses inverted residuals and linear bottlenecks to build an efficient deep feature extractor~\cite{sandler2018mobilenetv2}. EfficientNet-B0 compounds depth, width, and resolution scaling for strong ImageNet performance with fewer parameters~\cite{tan2019efficientnet}. Yosinski et al.\ show that lower CNN layers learn generic features (edges, textures) transferable across domains, while upper layers are task-specific and benefit from fine-tuning~\cite{yosinski2014transferable}; this motivates our OFAT ablation across freezing depths.

Probabilistic calibration is especially important for selective classification~\cite{niculescu2005predicting}. Guo et al.\ show that modern neural networks are often overconfident and that temperature scaling is a simple, effective post-hoc fix~\cite{guo2017calibration}. We apply scalar temperature scaling~\cite{guo2017calibration} fitted on training OOF predictions to avoid any validation leakage.

Late fusion combines independent class-probability vectors from complementary models, reducing variance when their errors are largely disjoint~\cite{wolpert1992stacked}. Razavian et al.\ demonstrate that CNN features off-the-shelf outperform many purpose-built representations~\cite{razavian2014cnn}; our work extends this by combining them with domain-engineered features rather than replacing them.

%% ======================================================
\section{Data and Regularisation}
%% ======================================================
% Author: Om Yadav
\label{sec:data}
\textit{(Data pipeline and regularisation: Om Yadav.)}

We reuse the Phase~1 audit: 2,521 images partitioned group-aware into 1,764 training, 380 validation, and 377 test images (seed 42, manifest fingerprint \texttt{affd3e9d}). No new images are added in Phase~2 and the outer split is immutable.

For deep learning, images are resized to $224\times224$ with letterbox padding (white fill to preserve aspect ratio). Training augmentation applies random resized crop (scale $0.8$--$1.0$), horizontal flip ($p=0.5$), rotation ($\pm10^\circ$), and color jitter (brightness/contrast 0.15, saturation 0.10), followed by ImageNet mean/std normalization. Validation and test images use center crop only. Balanced sampling and class-weighted loss are ablated as alternatives to standard training (see Sec.~\ref{sec:ablation}).

Augmentation is implemented as a togglable \texttt{build\_transforms(train=True/False)} call that applies strictly to training loaders; validation and test use the identical deterministic pipeline. Five frozen training folds from Phase~1 are reused for all cross-validation, ensuring zero data leakage between selection and final evaluation.

%% ======================================================
\section{Architectures}
%% ======================================================
% Author: Aditya Bharadwaj
\label{sec:arch}
\textit{(Backbone architectures: Aditya Bharadwaj.)}

Two pretrained ImageNet backbones are evaluated:
\begin{itemize}
\item \textbf{MobileNetV2}~\cite{sandler2018mobilenetv2}: 2.2M parameters, inverted residual blocks with depthwise separable convolutions, 1,280-dimensional penultimate embedding. We index its 19 structural layers (features 0--18); ``last 2 blocks'' refers to features 17--18 (the final InvertedResidual and the $1\times1$ pointwise conv).
\item \textbf{EfficientNet-B0}~\cite{tan2019efficientnet}: compound-scaled variant with squeeze-and-excitation, 1,280-dimensional embedding.
\end{itemize}
In the frozen baseline (\textsc{Base}), all backbone weights are frozen and only a 6-class linear head with $p=0.2$ dropout is trained (7,686 parameters). In configuration \texttt{A2.1}, the last two MobileNetV2 blocks are unfrozen and trained at learning rate $10^{-4}$, while the head uses $10^{-3}$; total trainable parameters: 893,766.

Deterministic cached embeddings (SHA-256 provenance, stale-cache rejection) are pre-extracted for the H1 and H2 hybrid candidates evaluated in Sec.~\ref{sec:hybrid}.

%% ======================================================
\section{Training Procedure}
%% ======================================================
% Author: Shourya Bafna
\label{sec:training}
\textit{(Training loop and scheduling: Shourya Bafna.)}

Training uses AdamW~\cite{loshchilov2019adamw} with decoupled weight decay $10^{-4}$, cosine annealing with warm restarts~\cite{loshchilov2017sgdr} and $T_{\max}=10$. Frozen backbone blocks are kept in \texttt{eval()} mode (BatchNorm and stochastic depth off) throughout training. Early stopping triggers when validation fold loss fails to improve for five consecutive epochs; maximum budget is 40 epochs.

The final A2.1 model is trained for exactly 13 epochs on the full 1,764 training images (the median best epoch from 5-fold CV), without touching validation data, and saved as \texttt{artifacts/phase2/final/A2.1/model.pt} (SHA-256 \texttt{26592f8b...5d28e}). Out-of-fold probability vectors are exported for downstream late fusion (H3) and temperature calibration.

%% ======================================================
\section{Ablation Study}
%% ======================================================
% Author: Aditya Bharadwaj
\label{sec:ablation}
\subsection{Deep Learning Ablation Study}
\label{sec:dl_ablations}
To isolate the empirical contributions of backbone architecture, fine-tuning depth, image resolution, head regularisation, data augmentation components, and class imbalance mitigation, we conducted a pre-registered one-factor-at-a-time (OFAT) ablation study comprising 13 configurations across the 5 frozen training cross-validation folds (65 evaluated folds, evaluated on Apple Silicon GPU hardware).
The baseline model---a frozen MobileNetV2~\cite{sandler2018mobilenetv2} backbone with a linear classification head, $p=0.2$ dropout, and comprehensive data augmentation at $224 \times 224$ resolution---achieved a mean Macro~F1 of $0.8515 \pm 0.0235$ (87.19\% accuracy).
Under our pre-registered selection rule ($|\Delta| > \sigma = 0.0235$), swapping the frozen backbone to EfficientNet-B0~\cite{tan2019efficientnet} ($0.8557 \pm 0.0263$), adjusting dropout ($p=0.0$: $0.8503$, $p=0.5$: $0.8479$), omitting individual augmentations (e.g., removing $\pm 10^\circ$ rotation: $0.8490$), removing all augmentation ($0.8339 \pm 0.0204$, $\Delta = -0.0176$), and incorporating class weighting ($0.8453$) or balanced sampling ($0.8419$) yielded score differences within $\sigma$, exhibiting no clear difference under the pre-registered rule in overall Macro~F1.
However, class-weighted loss and balanced sampling explicitly improved mean per-fold trash recall from 65.63\% to 72.89\% and 72.84\% respectively (pooled out-of-fold recall from 63/96 to 70/96).
Conversely, reducing input resolution to $160 \times 160$ fell just beyond $\sigma$ (worse, Macro~F1 $0.8271 \pm 0.0257$, $\Delta = -0.0244$).
Crucially, partial fine-tuning of only the final two structural blocks in MobileNetV2 (\texttt{A2.1}, training 893,766 parameters while keeping the first 17 blocks frozen, with differential learning rates of $10^{-3}$ on the head and $10^{-4}$ on the backbone) produced a clear improvement under the rule to \textbf{0.8825 $\pm$ 0.0219 Macro~F1} (89.80\% accuracy, $\Delta = +0.0310 > \sigma$, median best epoch 13), driven by recall gains on trash (79.21\% vs.\ 65.63\%) and metal (94.39\% vs.\ 89.50\%).
We note as a limitation that most frozen configurations reached their best epoch at 38--40 of the pre-registered 40-epoch cap (e.g., \textsc{Base} median best 38, \texttt{A4.2} and \texttt{A5.1} median best 40) whereas \texttt{A2.1} early-stopped naturally (mean 20.2 epochs run); part of \texttt{A2.1}'s margin may reflect this budget constraint, and extended frozen training is left as future work.
Configuration \texttt{A2.1} is on par with the Phase 1 SVM on handcrafted features ($0.8803 \pm 0.0215$, $\Delta = 0.0022 \ll \sigma$), while Daksh's early feature fusion hybrid H2 ($0.9044 \pm 0.0332$) shows no clear gain over \texttt{A2.1} under the $\sigma$ rule ($\Delta = 0.0219 < 0.0332$), with the final evaluation deferred to upcoming late fusion (H3) and the held-out test set.
Consequently, configuration \texttt{A2.1} was designated as the primary deep-learning model, and its out-of-fold predictions were exported for late fusion.



%% ======================================================
\section{Hybrid Models}
%% ======================================================
% Author: Daksh Batra
\label{sec:hybrid}
% Phase 2 Hybrid ML + DL paragraph for IEEE report
% Author: Daksh Batra

\subsection{Hybrid ML and Deep Learning Feature Fusion}
While handcrafted photometric and texture descriptors achieved strong baseline classification in Phase~1 ($0.8803\pm0.0215$ mean cross-validation macro F1 across five folds), systematic confusions persisted among glass, metal, and plastic due to overlapping specular reflections and smooth surface gradients. To investigate whether high-level geometric representations resolve these material confusions, we evaluated a hybrid architecture combining off-the-shelf deep convolutional neural network (CNN) representations~\cite{razavian2014cnn} with engineered descriptors. Pretrained 1,280-dimensional feature embeddings from frozen MobileNetV2 and EfficientNet-B0 backbones were extracted and fused with Phase~1's winning 1,608-dimensional extended color and spatial texture features. To prevent the higher-dimensional handcrafted block from dominating the radial basis function (RBF) kernel distance metric, we enforced a fold-local dimensional normalization protocol: each block was independently standardized using training-fold statistics and scaled by $1/\sqrt{D_{\mathrm{block}}}$ prior to concatenation ($D_{\mathrm{total}} = 2,888$), giving both modalities equal total geometric variance. Evaluating across the five frozen training cross-validation folds, early feature fusion with an RBF SVM ($C=10.0$, $\gamma=\text{`scale'}$) achieved a mean cross-validation macro F1 of $0.9044\pm0.0332$ and an overall accuracy of 91.73\% (compared to $0.8597\pm0.0194$ for classical classifiers trained on CNN embeddings alone). 

In accordance with our pre-registered decision rule ($\Delta > \sigma$), the $+0.0241$ absolute gain of early fusion under initial discrete predictions ($0.9044\pm0.0332$ Macro F1) did not constitute a clear win under the pre-registered rule as it remained within the cross-fold sample standard deviation ($\sigma = 0.0332$); when evaluated under the argmax of predicted probabilities for strict cross-model consistency, H2 achieves $0.9175\pm0.0262$ with 92.52\% pooled accuracy (1,632/1,764). Following the selection of the fine-tuned MobileNetV2 architecture (A2.1, $0.8825\pm0.0219$ macro F1, 89.80\% pooled accuracy) in Phase~2 ablations, we evaluated late probability fusion (H3) combining out-of-fold probability vectors from A2.1 and the calibrated Phase~1 SVM via convex blending ($p = w \cdot p_{\mathrm{cnn}} + (1 - w) \cdot p_{\mathrm{svm}}$). Analysis of out-of-fold predictions reveals that while CNN and SVM predictions disagree on 16.72\% of training images (295/1,764), both models fail simultaneously on only 2.61\% (46/1,764), demonstrating complementary error patterns. Using nested cross-validation across the five frozen training folds—where fusion weight $w \in [0, 1]$ was selected on inner training folds independently for each outer fold—equal weighting ($w=0.5$) was selected across all folds. H3 achieved a mean cross-validation macro F1 of $\mathbf{0.9364\pm0.0171}$, mean per-fold accuracy of 94.62\%, and pooled accuracy of $\mathbf{94.61\%}$ (1,669/1,764). Comparing against the best standalone single model (A2.1, 0.8825 macro F1), this represents a clear win under the pre-registered rule ($\Delta = 0.0539 > \sigma = 0.0219$). When compared to H2, $\Delta = 0.0189 < \sigma_{\mathrm{H2}} = 0.0262$, so H3 is carried forward as the highest mean candidate per protocol rather than a clearly separated winner under the rule over H2. All cross-validation estimates remain somewhat optimistic as base models were selected on these five folds; the final test split ($N=377$) provides the unbiased assessment. Pooled confusion analysis across all 1,764 training samples demonstrated that metal recall reached 98.25\% (281/286), reducing metal-to-glass errors to just 3 (an 85.7\% reduction compared to Phase~1), while plastic recall reached 91.67\% (308/336) and trash recall rose to 86.46\% (83/96). The final hybrid model was assembled on the full training split—preserving the verified A2.1 checkpoint (SHA-256: \texttt{26592f8b...5d28e}) without retraining—and exported for downstream probability calibration and operational abstention threshold selection on the held-out validation split.


%% ======================================================
\section{Calibration and Final Evaluation}
%% ======================================================
% Author: Prajjwal Tripathi
\label{sec:eval}
% Phase 2 Final Test Evaluation section for IEEE report
% Author: Prajjwal Tripathi (generated from saved evaluation JSONs)

\subsection{Frozen Phase 2 Test Evaluation and Statistical Comparison}
\label{sec:phase2_eval}
Before accessing the untouched TrashNet test split ($N=377$), all three finalists---Phase~1 RBF SVM, fine-tuned MobileNetV2 (\texttt{A2.1}), and equal-weight late fusion (\texttt{H3}, $w=0.5$)---were frozen in \texttt{protocol.json} (git revision \texttt{62db0847}). Each model's scalar temperature $T$ was fitted exclusively on its 1,764 training out-of-fold probabilities, and two abstention thresholds were selected strictly on the 380 validation images: a headline rule ($\ge 90\%$ accuracy, $\ge 80\%$ coverage) and a pre-declared demo rule ($\ge 95\%$ accuracy, $\ge 80\%$ coverage; added because \texttt{H3}'s headline threshold is $0.00$). Model SHA-256 hashes were verified prior to the single test pass.

\begin{table*}[t]
\caption{Untouched Phase~2 Test Results ($N=377$, 95\% Bootstrap CIs over 2,000 Resamples, and Both Pre-Declared Frozen Thresholds)}}
\label{tab:phase2_main}\centering
\begin{tabular}{lcccccc}\toprule
Model & $T$ & Macro F1 [95\% CI] & Accuracy [95\% CI] & ECE & Headline ($\tau$, Cov/Acc) & Demo ($\tau$, Cov/Acc)\\\midrule
Phase 1 SVM & 0.834 & 0.8963 [0.860, 0.926] & 90.19\% [87.0, 93.1] & 0.0361 & 0.45 (97.6\% / 91.6\%) & 0.76 (84.4\% / 97.2\%)\\
A2.1 CNN & 1.103 & 0.8577 [0.813, 0.897] & 88.06\% [84.6, 91.2] & 0.0367 & 0.50 (96.6\% / 90.4\%) & 0.87 (76.4\% / 95.5\%)\\
H3 Hybrid & 0.536 & \textbf{0.9233} [0.889, 0.952] & \textbf{93.63\%} [91.0, 96.0] & 0.0130 & 0.00 (100.0\% / 93.6\%) & 0.53 (96.6\% / 95.3\%)\\
\bottomrule
\end{tabular}
\end{table*}

Table~\ref{tab:phase2_main} reports the single test evaluation. \texttt{H3} Hybrid achieves the strongest test performance across every metric, reaching Macro~F1 $\mathbf{0.9233}$ and accuracy $\mathbf{93.63\%}$ with the lowest expected calibration error ($\mathrm{ECE}=0.0130$, Brier $0.0971$, log loss $0.2257$). At its frozen demo threshold ($\tau=0.53$), \texttt{H3} retains 96.55\% coverage at 95.33\% accepted accuracy, whereas \texttt{A2.1}'s strict demo threshold ($\tau=0.87$) drops to 76.39\% coverage on test. Compared against the historical Phase~1 test ensemble (0.8999 Macro~F1, 90.19\% accuracy, ECE 0.1325), the Phase~2 full-train \texttt{phase1\_svm} matches accuracy exactly (90.19\%, Macro~F1 0.8963) while OOF temperature scaling reduces test ECE by 72.8\% to 0.0361.

\begin{table}[t]
\caption{Pairwise Exact McNemar Tests on Unconditional Test Predictions ($\alpha_{\mathrm{Bonf}} = 0.05/3 = 0.0167$)}
\label{tab:phase2_mcnemar}\centering
\begin{tabular}{lrrrrrc}\toprule
Comparison (A vs.\ B) & Both OK & A Only & B Only & Both Err & $p$ & Bonf.\\\midrule
H3 Hybrid vs.\ A2.1 CNN & 330 & 23 & 2 & 22 & 0.0000 & Yes\\
H3 Hybrid vs.\ Phase 1 SVM & 330 & 23 & 10 & 14 & 0.0351 & No\\
A2.1 CNN vs.\ Phase 1 SVM & 307 & 25 & 33 & 12 & 0.3581 & No\\
\bottomrule
\end{tabular}
\end{table}

Pairwise exact McNemar binomial tests on discordant unconditional predictions (Table~\ref{tab:phase2_mcnemar}) confirm that \texttt{H3} Hybrid significantly outperforms \texttt{A2.1} CNN ($p=0.00002 < 0.0167$, 23 vs.\ 2 discordant wins). Although \texttt{H3} beats Phase~1 SVM nominally (23 vs.\ 10, $p=0.0351 < 0.05$), that margin does not clear the Bonferroni threshold $0.0167$ on $N=377$ images; the two single models do not differ significantly ($p=0.3581$). Hook-based Grad-CAM on \texttt{A2.1}'s final convolutional block across 12 test images (6 correct, 6 errors) gives a mixed test of the ``CNN sees shape'' hypothesis: 11/12 heatmaps place $\ge 80\%$ of activation mass on the foreground object (mean 93.6\%) rather than TrashNet's white posterboard, yet all six inspected errors involve clear object localization on overlapping glass, metal, and plastic shapes and specular reflections. External robustness remains unverified (robustness set not available on \texttt{main}), leaving TrashNet background dependence, CV selection optimism, and low test support ($n=20$ for trash) as primary limitations.


\subsection{Grad-CAM Interpretability}
\label{sec:gradcam}
Hook-based Grad-CAM~\cite{selvaraju2017gradcam} is applied to \texttt{A2.1}'s final convolutional block across a 12-image grid (6 correct spanning all classes, 6 high-confidence test errors spanning glass/metal/plastic confusions). The ``foreground'' heuristic is HSV saturation $> 0.12$ or value $< 0.92$ (a colour threshold, not an object mask); Table~\ref{tab:gradcam} reports the foreground activation share for each inspected case.


\begin{table}[t]
\caption{A2.1 Grad-CAM Foreground Activation Share}
\label{tab:gradcam}
\centering
\footnotesize
\setlength{\tabcolsep}{4pt}
\begin{tabular}{lllr}
\toprule
Image & Outcome & True$\to$Pred & FG CAM\\
\midrule
\texttt{cardboard1} & Correct & card$\to$card & 98.3\%\\
\texttt{glass10} & Correct & glass$\to$glass & 69.3\%\\
\texttt{metal103} & Correct & metal$\to$metal & 98.7\%\\
\texttt{paper106} & Correct & paper$\to$paper & 88.1\%\\
\texttt{plastic1} & Correct & plastic$\to$plastic & 86.9\%\\
\texttt{trash109} & Correct & trash$\to$trash & 98.0\%\\
\texttt{glass307} & Wrong & glass$\to$metal & 93.0\%\\
\texttt{glass102} & Wrong & glass$\to$plastic & 97.7\%\\
\texttt{plastic233} & Wrong & plastic$\to$glass & 97.1\%\\
\texttt{metal293} & Wrong & metal$\to$glass & 99.2\%\\
\texttt{plastic467} & Wrong & plastic$\to$metal & 97.2\%\\
\texttt{plastic398} & Wrong & plastic$\to$glass & 99.3\%\\
\bottomrule
\end{tabular}
\end{table}

\noindent\textbf{Key finding:} 11 of 12 heatmaps concentrate ${\geq}80\%$ of activation mass on the foreground object (mean 93.6\%), suggesting foreground-focused attention in these inspected examples, without ruling out background dependence. However, Grad-CAM \emph{does not} explain the performance gap between A2.1 and Phase~1 SVM (McNemar $p=0.358$, not significant); nor does it explain H3's gain, which is demonstrated by discordant-count analysis (23 vs 2, $p=0.000019$). Shape-focused attention is insufficient for glass/metal/plastic disambiguation because those materials share overlapping silhouette geometry.

%% ======================================================
\section{Demo}
%% ======================================================
\label{sec:demo}
The Gradio 6.x application loads all three Phase~2 models once per process from frozen artifacts and reads temperatures and demo thresholds from \texttt{protocol.json}. A model selector (H3 hybrid recommended, A2.1 CNN, Phase~1 SVM) lets users compare all three in a single interface. Predictions below the frozen demo threshold trigger an amber ``Uncertain --- manual inspection recommended'' banner. SHA-256 hash verification at startup rejects mismatched or missing model files. The demo is for educational purposes; local recycling rules determine actual disposal.

%% ======================================================
\section{Limitations}
%% ======================================================
\label{sec:limits}
\begin{enumerate}
\item \textbf{Cross-validation optimism:} All models reuse the same 5 training folds for model selection; reported CV scores are optimistic ($H3$ CV 0.9364 vs test 0.9233; A2.1 CV 0.8825 vs test 0.8577).
\item \textbf{Small test set ($N=377$, trash $n=20$):} With only 33 discordant predictions between H3 and SVM, the nominal $p=0.035$ does not survive Bonferroni correction ($\alpha_{Bonf}=0.0167$).
\item \textbf{TrashNet white background:} All images were photographed against a uniform white posterboard; the robustness set was not collected (\texttt{docs/robustness/} absent on \texttt{main}).
\item \textbf{OOF temperature applied to full-train models:} Temperature was fitted on fold-model OOF predictions but applied to the full-training final models, following Phase~1 practice.
\item \textbf{H3 from uncommitted working tree:} \texttt{protocol.json} records \texttt{working\_tree\_dirty} in H3's metadata; model and OOF SHA-256 hashes are pinned and verified before every test prediction.
\item \textbf{Threshold fragility:} A2.1's demo threshold (0.87) achieved 82.4\% validation coverage but only 76.4\% test coverage, showing that tight confidence thresholds do not transfer exactly across splits.
\end{enumerate}

%% ======================================================
\section{Conclusion}
%% ======================================================
\label{sec:conc}
Phase~2 shows that partial fine-tuning of the last two MobileNetV2 blocks matches the Phase~1 handcrafted SVM (A2.1 CV 0.8825 vs SVM 0.8803) and that late probability fusion of complementary error patterns yields a clear improvement (H3 test macro F1 \textbf{0.9233}, significantly better than A2.1, $p=0.000019$, Bonferroni-corrected). Temperature calibration on OOF data reduces ECE from 0.095 to 0.019 for H3. The frozen evaluation protocol with pre-registered thresholds ensures an auditable, single-pass test result.

\balance
\bibliographystyle{IEEEtran}
\bibliography{references}

%% Contributions table
\begin{table}[t]\caption{Phase~2 Technical Contributions}\centering
\begin{tabular}{lp{5.2cm}}\toprule
Member & Primary contribution\\\midrule
Om Yadav & Phase~2 data pipeline, augmentation, balanced sampling\\
Aditya Bharadwaj & Backbone loaders, OFAT ablation study (13 configs, 65 folds)\\
Shourya Bafna & Training loop, AdamW/cosine schedule, early stopping\\
Daksh Batra & H1/H2 hybrids, H3 late fusion, error complementarity analysis\\
Prajjwal Tripathi & Temperature scaling, abstention thresholds, frozen test eval\\
Saksham Miglani & Demo, Phase~2 report and presentation assembly\\\bottomrule
\end{tabular}\end{table}

\end{document}
